⟦BLK-0⟧
⟦CAP-0⟧ Scaling Behaviour of Section Aware Chunking
⟦CAP-1⟧ A study of how translation units grow when the section tree is respected.

\section{Introduction}
The introduction is deliberately short so that it can be aggregated with the section
that follows it whenever the soft target allows such aggregation to happen.

\section{Related work}
Related work is also short, which makes it a good candidate for aggregation with the
introduction under a generous soft target for translation units.

\section{Method}
The method section is long enough that it has to be split, and it carries subsections
so that the splitter can use structural boundaries instead of raw paragraph counts.

\subsection{Masking}
Masking walks the token stream and replaces every heavy environment with an opaque
placeholder, so the translator never sees a table body or a display equation at all.
The classification table decides which environments count as heavy in this pass.

Unknown environments are masked conservatively, which costs coverage but never costs
correctness, and the round trip self check proves the decision on every single paper
before a single token of budget has been spent on the translation itself.

\subsection{Chunking}
Chunking prefers whole sections because coherence inside a section is where most of
the translation quality comes from, while coherence between sections is weak anyway.
Terminology consistency across sections is the job of the glossary and the brief.

The size constraint comes from the output side rather than the input side, since the
translation is roughly as long as the original and very long single generations drift,
get lazy towards the end, and eventually run into the output limit of the model.

\subsection{Validation}
Validation compares placeholder multisets, control sequence multisets, brace balance,
and paragraph counts, so any block can be split back into aligned paragraph pairs.
That alignment is what keeps the fallback granularity at the paragraph level.

Because the fallback unit stays small, making the translation unit large costs nothing
in terms of degradation, which is the whole argument for section first chunking as a
default strategy rather than a fixed window over characters.

\section{Experiments}
The experiment section has no subsections at all, so an oversized experiment section
must fall back to paragraph boundaries when it is divided into translation units.

We compare three chunking strategies on a corpus of papers drawn from several fields.

The first strategy is a fixed character window, which is the baseline inherited from
the prototype and which ignores the structure of the document completely.

The second strategy respects paragraph boundaries but still ignores the section tree.

The third strategy is the one described here, and it wins on both retry rate and on
terminology consistency across the whole document.

Results are reported as averages over five runs with a fixed seed for the sampler.

\section{Notation}
This section is a single very long paragraph on purpose, because a paragraph is the
atom of the chunker and must never be divided, no matter how large it grows in the
source document, and the placeholder ⟦BLK-1⟧ inside it does not change that rule at
all, so the resulting translation unit is allowed to exceed the hard limit while the
splitter reports it as one indivisible paragraph, which is exactly the behaviour that
the paragraph alignment in the validation stage relies upon later in the pipeline,
and which the compile stage relies upon again when it bisects a failing block down to
the individual paragraph that broke the build, restores the original text for it, and
keeps every other paragraph of the same block in its translated form regardless.

\section{Conclusion}
Short.
